\section{模型假设}
为简化问题、突出主要物理过程，本文做出以下假设：
\begin{itemize}
	\item 假设 1（几何简化）：药材近似为\textbf{无限长圆柱}，因长径比 $25/4=6.25$ 较大，忽略端面传热传质，温度与水分浓度仅沿径向变化，即一维径向模型。
	\item 假设 2（物性均匀）：药材初始温度与含水率均匀分布，物性各向同性；导热系数、比热容、密度与水分扩散系数仅随含水率（及温度）按题目给定的经验公式变化。
	\item 假设 3（边界条件）：药材表面与烘房空气之间的对流换热、对流传质分别由牛顿冷却定律与对流传质定律描述，对流换热系数 $h$、对流传质系数 $h_m$ 为常数（问题 2--4 沿用附录 2 的 $h=25$ W/(m$^2\cdot$K)、$h_m=8\times10^{-7}$ m/s）。
	\item 假设 4（烘房环境）：预热段（$0\sim4$ h）烘房温度与空气水分浓度按附件 1 实测值线性插值；恒温干燥段（$4$ h 之后）取恒定值 $T_{air}=50$ °C、$C_{air}=0.050$ kg/kg。
	\item 假设 5（收缩处理，问题 4）：药材收缩时各层按比例均匀收缩，半径 $R(t)$ 由附件 2 数据经保单调三次样条（PCHIP）插值给出；在固定网格上采用移动边界近似，忽略收缩引起的对流输运项（中心处材料速度为零，其对干燥终点影响有限）。
	\item 假设 6：干燥过程中无化学反应与内部相变，水分以扩散方式迁移，忽略重力与辐射换热的影响。
\end{itemize}

\section{符号说明}
本文使用的主要符号及其含义如表~\ref{tab:symbols} 所示。

\begin{table}[htbp]
	\caption{符号说明}
	\label{tab:symbols}
	\centering
	\begin{tabular}{clc}
		\toprule
		符号 & 说明 & 单位 \\
		\midrule
		$r$ & 到药材中心的径向距离 & cm \\
		$R$ & 药材半径（随时间变化） & cm \\
		$t$ & 时间 & s / h \\
		$T(r,t)$ & 药材温度 & °C \\
		$C(r,t)$ & 干基含水率（水分浓度） & kg/kg \\
		$T_{air}(t)$ & 烘房空气温度 & °C \\
		$C_{air}(t)$ & 烘房空气水分浓度 & kg/kg \\
		$\rho$ & 密度 & kg/m$^3$ \\
		$c_p$ & 比热容 & J/(kg$\cdot$K) \\
		$k$ & 热传导系数 & W/(m$\cdot$K) \\
		$D$ & 水分浓度扩散系数 & m$^2$/s \\
		$h$ & 对流换热系数 & W/(m$^2\cdot$K) \\
		$h_m$ & 对流传质系数 & m/s \\
		$T_0$ & 初始温度（$28$ °C） & °C \\
		$C_0$ & 初始含水率（$2.55$ kg/kg） & kg/kg \\
		$\Delta r,\Delta t$ & 空间、时间步长 & m, s \\
		\bottomrule
	\end{tabular}
\end{table}
